\section{Data-Analysis Agents and Known-Answer Evaluation}\label{sec:data}

Data analysis is the part of automated research that is easiest to grade, because many analytic questions have one right answer that can be computed before the agent runs. The benchmark literature supports a two-part conclusion. Accuracy against known answers is moderate on closed-form questions over a single table and low on anything multi-step, and the published numbers move in both directions for reasons that have little to do with analytic skill. On InfiAgent-DABench, which grades closed-form questions over single CSV files by exact match, GPT-4 reached 78.99\% \citep{D1}. On harder, multi-step tasks, the best systems evaluated in 2024 scored 30.5\% on DA-Code, 34.12\% on DSBench's analysis tasks and 25\% on DiscoveryBench, and GPT-4 reached 58\% on QRData's statistical and causal questions, most of which use a single data sheet \citep{D2,D3,D6,D14}. Within one benchmark, difficulty dominates: on DABstep, o4-mini answered 76.39\% of easy tasks and 14.55\% of hard ones \citep{D9}. Newer models score higher without scoring high; on a multi-database benchmark the best agent reached 57\% pass@1 \citep{D11}. Tuning to a benchmark inflates scores. On that benchmark's live, self-submitted leaderboard (\repofile{ucbepic/DataAgentBench}{README.md}, read 21 November 2026), the top row reports 0.9917 pass@1 with a benchmark-specific prompt and a note for each query that ``tells the agent how to solve that query''; the best row without a tuned prompt reports 0.8858. The leaderboard uses a different query set from the paper (54 queries against 104), so neither figure is comparable with the paper's 57\%, and the top rows partly measure tuning to the benchmark rather than general ability.

\subsection{Plausible Wrong Numbers}

The dominant failure is not a crash but a well-formed, wrong answer. On the multi-database benchmark, wrong plans and wrong implementations account for 90\% of incorrect answers; in one example, a regular expression for \textit{MALE} matches inside \textit{FEMALE} and mislabels gender \citep{D11}. In QRData, 44\% of GPT-4's sampled errors were knowledge-recall failures ``like choosing a wrong hypothesis test, or using statistical methods to address a causal question'' \citep{D14}. The strongest evidence is an audit, by researchers outside the developer, of DS-STAR, a Google Cloud and KAIST agent that already checks its plan with an LLM verifier \citep{D16}. Of 580 zero-scored DS-STAR runs, 405 (69.8\%) ``delivered normal-looking but incorrect results without reporting failure'' \citep{D12}. That is one preprint on one agent, but it fits a second result: on DiscoveryBench, planning with a self-critic did not improve a simple code-generation agent, while oracle feedback did \citep{D6}. The gain comes from a signal outside the agent.

\subsection{The Right Conclusion, the Wrong Number}

Statistics-specific benchmarks separate the verdict from the estimate, and the two diverge. On P-Bench's hard split, GPT-5.4 reached the correct reject or fail-to-reject decision in 58.3\% of tasks but a p-value close to the canonical analysis in only 30.5\% \citep{D21}. Those authors also built an agent that outperforms GPT-5.4 on their benchmark, but the gap reported here is measured on GPT-5.4, a model they did not build. CausalReasoningBenchmark finds the same split for one baseline LLM: the high-level identification strategy was right in 81\% of cases and the full specification in 38\% \citep{D20}. A grader that checks only the verdict overstates reliability, a pattern Section~\ref{sec:replication} finds in human and agent analysts alike.

\subsection{Imperfect Data and Imperfect Keys}

Success on clean data does not carry over. On RADAR, o4-mini scored 100\% on clean tables with code execution, and its performance dropped by 59\% on perturbed ones \citep{D17}. An audit of agent-written modelling pipelines found that every framework--model configuration had at least one failure mode in 43\% or more of its pipelines, and that every data-leakage mode improved a pipeline's ranking \citep{D23}; scoring by performance alone rewards leakage, although that audit used smaller models (two Qwen models that run on commodity hardware and Gemini 3.7 Flash) and detected failures with an LLM judge. A vendor article by Anthropic, on its own internal evaluations, reports accuracy that ``didn't exceed 21\%'' without curated skills, rose ``consistently above 95\%'' with them, then drifted ``from ${\sim}$95\% at launch to ${\sim}$65\% over a month'' as the skill documents went unmaintained; ``the failure mode none of this fully catches is the silent one'' \citep{D26}.

The answer key is part of the instrument, and it is wrong often enough to matter. An audit of a data-engineering benchmark found that 82.7\% of 81 failed tasks contained at least one benchmark-attributable error; fixing the key raised the same agent and model from 22.66\% to 32.51\% \citep{D13}. An automated auditor found 12 defects in ScienceAgentBench, all confirmed by its authors, despite ``multiple rounds of manual validation'' \citep{D15}. DABench kept a label only when OpenAI's ADA agent gave the same answer in all three trials, before human filtering \citep{D1}, so its answer key rests on one agent's agreement with itself. Two open issues report apparently wrong labels (\repo{InfiAgent/InfiAgent}{28}, \repo{InfiAgent/InfiAgent}{29}). Reading its grader (\repofile{InfiAgent/InfiAgent}{examples/DA-Agent/eval_closed_form.py}) shows that it scores only questions with a non-empty response and divides by the number scored, so a run that skips hard questions scores higher.

Known-answer evaluation is the strongest external signal available for analysis agents, and the evidence suggests it holds only under conditions that are easy to omit: an audited key, the full task set as the denominator, a grade on the estimate as well as the verdict, and tests on perturbed data and for leakage.
